\section{Introduction}

%%%%%%%%%%%%%%%%%%%%%%%%%%%%%%%%%%%%%%%%

\begin{comment}
\begin{frame}{News}
  \centering

  %  Ajuste: tira 1 cm de cada lado ─ mexa nesses números até “encher” de fato
  \begin{adjustwidth*}{-1.5cm}{-1.5cm}

  % Bloco horizontal com largura exata da lâmina
  \makebox[\paperwidth][c]{%
    \includegraphics[width=0.32\paperwidth,
                     page=1,trim=0 5 0 15,clip]{Noticias extremas/noticia_1.pdf}%
    \hspace{0.008\paperwidth}%
    \includegraphics[width=0.32\paperwidth,
                     page=1,trim=0 5 0 15,clip]{Noticias extremas/noticia_2.pdf}%
    \hspace{0.008\paperwidth}%
    \includegraphics[width=0.32\paperwidth,
                     page=1,trim=0 5 0 15,clip]{Noticias extremas/noticia_3.pdf}%
  }

  \vspace{-0.3cm}  % pequeno espaço entre imagem e legendas

  % Legendas alinhadas horizontalmente sob cada imagem
  {\scriptsize
    \hspace*{0.1\paperwidth} Al Jazeera, 2025
    \hspace*{0.15\paperwidth} European Commission, 2025
    \hspace*{0.13\paperwidth} BBC News, 2025
  }
\end{adjustwidth*}
\end{frame}
\end{comment}

\begin{frame}{Context}
  \centering
  \includegraphics[
    width=\textwidth,
    height=0.85\textheight,
    keepaspectratio
  ]{Noticias extremas/news1.pdf}
\end{frame}

\begin{frame}{Context}
  \centering
  \includegraphics[
    width=\textwidth,
    height=0.85\textheight,
    keepaspectratio
  ]{Noticias extremas/news2.pdf}
\end{frame}

\begin{frame}{Extreme events}
  \centering
  \includegraphics[
    width=\textwidth,
    height=0.85\textheight,
    keepaspectratio
  ]{Noticias extremas/extremes.png}
\end{frame}


%%%%%%%%%%%%%%%%%%%%%%%%%%%%%%%%%%%%%%%%

\begin{comment}
%%%%%%%%%%%%%%%%%%%%%%%%%%%%%%%%%%%%%%%%
\begin{frame}{Motivation}
  \centering

  %  Ajuste: tira 1 cm de cada lado ─ mexa nesses números até “encher” de fato
  \begin{adjustwidth*}{-1.5cm}{-1.5cm}

  % Bloco horizontal com largura exata da lâmina
  \makebox[\paperwidth][c]{%
    \includegraphics[width=0.32\paperwidth,
                     page=1,trim=0 5 0 15,clip]{Noticias extremas/noticia_4.pdf}%
    \hspace{0.008\paperwidth}%
    \includegraphics[width=0.32\paperwidth,
                     page=1,trim=0 5 0 15,clip]{Noticias extremas/noticia_5.pdf}%
    \hspace{0.008\paperwidth}%
    \includegraphics[width=0.32\paperwidth,
                     page=1,trim=0 5 0 15,clip]{Noticias extremas/noticia_6.pdf}%
  }

  \vspace{-0.3cm}  % pequeno espaço entre imagem e legendas

  % Legendas alinhadas horizontalmente sob cada imagem
  {\scriptsize
    \hspace*{0.045\paperwidth} NASA Earth Observatory, 2025
    \hspace*{0.09\paperwidth} The New York Times, 2025
    \hspace*{0.14\paperwidth} SciTechDaily, 2025
  }
\end{adjustwidth*}
\end{frame}



%%%%%%%%%%%%%%%%%%%%%%%%%%%%%%%%%%%%%%%%
\begin{frame}{Historical context}

\begin{enumerate}
    \item Fisher and Tippett (1928): pioneering results in Extreme Value Theory (EVT)
        \vspace*{0.7cm}
    \item Von Mises (1936): fundamental contributions to the theoretical foundations of EVT
        \vspace*{0.7cm}        
    \item Gnedenko (1943): rigorous formalization of the limiting extreme value distributions
        \vspace*{0.7cm}
    \item Jenkinson (1955): establishment of the Generalized Extreme Value (GEV) distribution
        \vspace*{0.7cm}   
    \item 1940s–1955 to present: GEV as the primary choice for modeling univariate extremes
        \vspace*{0.7cm}   
    \item Recognized limitations of the GEV under complex or multimodal scenarios
\end{enumerate}


% refs
\begin{flushright}
{\tiny Referências: Haan e Ferreira, 2006; Davison, Padoan e Ribatet, 2012; Kotz e Nadarajah, 2000; Otiniano \textit{et al.}, 2023b}
\end{flushright}


\end{frame}


%%%%%%%%%%%%%%%%%%%%%%%%%%%%%%%%%%%%%%%%
\begin{frame}{Contexto histórico}
  \setbeamercovered{transparent}

  \begin{enumerate}
  \setcounter{enumi}{5}
    \item Surgem várias extensões e generalizações da GEV, por exemplo:\\
    \begin{itemize}
      \item Aryal e Tsokos (2009): Transmutada GEV (TGEV)
      \item Provost \textit{et al.} (2018): q-GEV
      \item Gramosa \textit{et al.} (2020): GEV inflada de zeros (ZIGEV)
      \item Otiniano \textit{et al.} (2023): GEV Bimodal (BGEV)
    \end{itemize}
  \end{enumerate}

  \pause
  \vspace*{0.5cm}

  A \textbf{BGEV} de Otiniano \textit{et al.} (2023):
  \begin{itemize}
    \item Tem um parâmetro a mais;
    \item Apresenta mais flexibilidade que a GEV; e
    \item Captura dados extremos com heterogeneidade e múltiplos regimes
  \end{itemize}

\end{frame}
\end{comment}


